Diplomová práce posuzuje stav a~účinnost současného sociálního dialogu v~České republice v~evropském srovnání a~vyhledává právní a~organizační nástroje přenositelné z~vybraných evropských systémů sociálního dialogu. Hypotéza práce zní: \emph{Posílení sociálního dialogu a~kolektivního vyjednávání v~České republice by přispělo k~udržitelnému hospodářskému rozvoji.} Je~posuzována třemi metodami: popisná analýza dokládá zhoršující se kondici sociálních partnerů (klesající odborová a~zaměstnavatelská organizovanost, nízké pokrytí kolektivními smlouvami, velmi nízké výdaje na~aktivní politiku zaměstnanosti), komparace se~systémy v~Dánsku, Rakousku, Německu, Polsku a~na~Slovensku zařazuje Českou republiku mezi ekonomiky s~dominantním podnikovým vyjednáváním a~korelační analýza ukazatelů trhu práce v~panelu zemí EU ukázala statisticky významnou vazbu mezi pokrytím kolektivními smlouvami, produktivitou práce a~průměrnou pracovní dobou. Analýza je doplněna čtyřmi případovými studiemi podnikových a~odvětvových smluv. Práce na~základě těchto zjištění formuluje návrhy inovací sociálního dialogu, mezi něž~patří aktivace institutu rozšiřování závaznosti odvětvových kolektivních smluv, převzetí části aktivní politiky zaměstnanosti odbory, mzdová transparentnost nad~rámec minimální transpozice směrnice EU 2023/970, vyrovnání postavení samostatně výdělečně činných a~platformových pracovníků v~kolektivním vyjednávání a~veřejný digitální registr kolektivních smluv.
